\documentclass[8pt]{extarticle}

\usepackage[margin=0.3in]{geometry}
\usepackage{multicol}
\usepackage{setspace}
\usepackage{titlesec}

% Compact spacing
\setstretch{0.9}
\setlength{\parskip}{0pt}
\setlength{\parindent}{0pt}

% Tight section spacing
\titlespacing{\section}{0pt}{1ex}{0.5ex}
\titlespacing{\subsection}{0pt}{0.8ex}{0.3ex}

\begin{document}

\begin{multicols}{3}
    \section{fMRI for Program Comprehension}
    Motivation: Indirect measures (performance, surveys) fail to capture actual cognition in code understanding. Methodology: fMRI used to observe brain activity while developers perform comprehension tasks with controlled baselines. Key Findings: Code comprehension activates left-hemisphere regions tied to working memory, attention, and language, indicating a hybrid cognitive process distinct from pure math or language. Implications: Enables objective measurement of cognitive load, allowing evaluation of languages, tools, expertise, and code complexity.

    \section{Code Readability and Cognitive Load}
    Motivation: Poor naming and readability are suspected to increase effort but lack strong empirical validation. Methodology: Controlled experiments with varied identifier quality and readability, measuring task accuracy, time, and physiological cognitive load. Key Findings: Low-quality identifiers and poor structure significantly increase mental effort; combined effects further degrade comprehension and performance. Implications: Readability is a core quality attribute; enforcing naming conventions and clean structure directly improves developer efficiency.

    \section{Code Review and Expertise in the Brain}
    Motivation: It is unclear whether code is processed like natural language or as a distinct cognitive task. Methodology: fMRI experiments comparing code comprehension, code review, and prose reading, with machine learning classification of neural patterns. Key Findings: Code tasks are distinguishable from prose; however, experts show more overlap, suggesting abstraction and familiarity. Implications: Programming is a unique cognitive domain, and expertise changes mental representations, impacting education and tool design.

    \section{Learning a New Programming Language}
    Motivation: Developers struggle with new languages despite prior experience due to transfer of incorrect assumptions. Methodology: Mixed-method study combining Stack Overflow analysis and developer interviews to identify learning patterns. Key Findings: Cross-language interference is common; developers rely on opportunistic, just-in-time learning, leading to shallow or incorrect mental models. Implications: Learning tools should explicitly address differences between languages and support conceptual transitions, not just syntax learning.

    \section{Predictors of Developer Productivity}
    Motivation: Productivity is often mismeasured, with unclear drivers across environments. Methodology: Large-scale survey with regression analysis evaluating technical, social, and organizational factors. Key Findings: Social factors (motivation, peer support, feedback) consistently outweigh technical factors; interruptions and environment strongly influence perceived productivity. Implications: Improving productivity requires better team culture, communication, and workflow design, not just better tools.

    \section{Developer Productivity and Code Quality at Google}
    Motivation: Prior work shows correlations but lacks causal evidence for productivity drivers. Methodology: Longitudinal analysis combining surveys and engineering logs using fixed-effects and lagged regression models. Key Findings: Code quality is the strongest predictor of future productivity; improvements in code health precede efficiency gains; tool satisfaction and stable priorities also matter. Implications: Investing in maintainability and reducing technical debt yields immediate productivity benefits.

    \section{Logs-Based Focus and Flow}
    Motivation: Flow is important but difficult to measure without intrusive self-reporting. Methodology: Multi-phase study using developer logs, embeddings of actions, and validation against surveys to derive a behavioral focus metric. Key Findings: Pure flow cannot be directly inferred, but focused work can be detected; frequency of focus sessions correlates strongly with perceived productivity. Implications: Provides scalable, passive measurement of focus, enabling organizations to study interruptions and optimize work patterns.

    \section{GenAI, Creativity, and Software Development}
    Motivation: Research focuses on productivity gains from AI, while its effect on creativity is underexplored. Methodology: Conceptual analysis using the 4P creativity framework and McLuhan’s tetrad to explore potential impacts. Key Findings: AI enhances ideation and reduces routine effort but risks homogenization, reduced originality, and over-reliance on generated solutions. Implications: Calls for research to ensure AI augments human creativity and supports diverse, novel problem-solving approaches.

    \section{Beliefs, Practices, and Personalities of Engineers}
    Motivation: The relationship between personality traits, engineering practices, and beliefs is not well understood. Methodology: Large-scale industrial survey using the Five-Factor Model along with questions on practices and preferences. Key Findings: Strong agreement on practices like testing and readability; disagreement on tools and environments; minimal personality differences across roles, with slight variation for managers. Implications: Engineering culture is broadly shared, and effectiveness depends more on context and practices than inherent personality differences.

    \section{Requirements Elicitation Survey}
    Motivation: Requirements elicitation is critical yet complex due to stakeholder diversity, ambiguity, and evolving needs. Methodology: Comprehensive survey categorizing elicitation techniques (interviews, questionnaires, observation, prototyping) and approaches across contexts. Key Findings: No single technique is sufficient; effectiveness depends on project type, stakeholder availability, and domain complexity; hybrid approaches are most common. Implications: Practitioners should select techniques contextually and combine methods; tool support and structured processes are essential to improve elicitation quality.

    \section{RE Techniques and Approaches Comparison}
    Motivation: Practitioners struggle to choose appropriate elicitation techniques due to lack of clear comparisons. Methodology: Comparative analysis of techniques based on cost, effectiveness, scalability, and stakeholder involvement. Key Findings: Interviews provide depth but are costly; surveys scale but lack richness; observation captures real behavior but is time-intensive; prototyping improves clarity but may bias requirements. Implications: Trade-offs must be considered explicitly; combining techniques mitigates individual weaknesses and improves requirement completeness.

    \section{Tool Support in Requirements Engineering}
    Motivation: Manual elicitation processes are error-prone and difficult to scale. Methodology: Survey of tools supporting documentation, collaboration, traceability, and automation in RE workflows. Key Findings: Tools improve organization and traceability but cannot replace human judgment; adoption depends on usability and integration with workflows. Implications: Effective RE requires both tool support and human-centered processes; future tools should better integrate automation (e.g., AI) with stakeholder interaction.

    \section{Issues and Pitfalls in Requirements Engineering}
    Motivation: RE failures often lead to project delays and defects, yet common pitfalls persist. Methodology: Analysis of recurring issues across studies and industry reports. Key Findings: Ambiguity, incomplete requirements, communication gaps, and stakeholder misalignment are dominant issues; over-reliance on a single technique worsens outcomes. Implications: Emphasizes need for iterative elicitation, validation, and stakeholder engagement throughout development.

    \section{Trends and Challenges in RE}
    Motivation: Software systems are becoming more complex, requiring evolution in RE practices. Methodology: Survey of emerging trends and open research challenges. Key Findings: Increasing importance of agile RE, continuous elicitation, and integration with AI-driven tools; challenges include handling dynamic requirements and large-scale stakeholder input. Implications: RE must become more adaptive, data-driven, and tool-supported to remain effective in modern development environments.

    \section{Conversational AI in RE Education}
    Motivation: AI tools like ChatGPT are increasingly used in education, but their role in RE learning is unclear. Methodology: Case study/course setup where students use conversational AI for RE tasks such as requirement drafting and refinement. Key Findings: AI assists in idea generation and structuring requirements but may introduce inaccuracies or superficial understanding; students rely on AI for efficiency. Implications: AI can enhance learning but must be guided with critical evaluation; educators should design tasks that balance assistance with deep understanding.

    \section{AI-Assisted Requirements Engineering}
    Motivation: AI has potential to automate parts of RE, but its effectiveness and risks are not fully understood. Methodology: Empirical and conceptual analysis of AI-supported RE tasks such as requirement extraction, summarization, and validation. Key Findings: AI improves speed and coverage but may introduce hallucinations, bias, and loss of context; human oversight remains essential. Implications: RE workflows should incorporate AI as an assistive tool rather than a replacement, with validation mechanisms to ensure correctness.

    \section{Human-AI Collaboration in Software Engineering}
    Motivation: AI tools are transforming development, but their broader impact on workflows and cognition is still emerging. Methodology: Synthesis of empirical studies and conceptual frameworks analyzing human-AI interaction in development tasks. Key Findings: AI improves productivity and exploration but affects reasoning, trust, and reliance patterns; developers shift from writing to supervising code. Implications: Future tools must support transparency, trust, and collaboration, ensuring developers remain engaged in reasoning rather than passive consumers.

    \section{AI, Productivity, and Developer Behavior}
    Motivation: Claims of productivity gains from AI tools require deeper understanding of behavioral changes. Methodology: Analysis of studies comparing developers using AI vs traditional workflows. Key Findings: AI increases speed and reduces effort for routine tasks, but may increase errors or reduce deep comprehension; novice vs expert usage differs significantly. Implications: Productivity gains must be balanced with quality and learning; AI tools should adapt to user expertise and task complexity.
\end{multicols}

\end{document}